Let $n_1<n_2<\cdots$ be integers with $n_i\ge 2$.
For a choice of residues $(a_i)_{i\ge 1}$ with $a_i\in\{0,\dots,n_i-1\}$, define for each $k\ge 1$ the uncovered set
\[U_k(a)=\{m\in\mathbb{Z}: m\not\equiv a_i\!\!\!\pmod{n_i}\ \text{for all }1\le i\le k\}.\]
This set is periodic modulo $L_k=\mathrm{lcm}(n_1,\dots,n_k)$, hence has a natural density $\mathbf{d}(U_k(a))$.
Define also the infinite uncovered set $U_\infty(a)=\bigcap_{k\ge 1}U_k(a)$; the hypothesis assumes $\mathbf{d}(U_\infty(a))=0$ for every choice $a$.

Question: Does $\mathbf{d}(U_\infty(a))=0$ for all residue choices $a$ imply the following uniform finite approximation property?
For every $\epsilon>0$ there exists $k$ such that for every $a$, $\mathbf{d}(U_k(a))<\epsilon$.
